\documentclass[10pt,a4paper]{article}

% ------------------------------------------------------------------ packages
\usepackage[T1]{fontenc}
\usepackage[utf8]{inputenc}
\usepackage{lmodern}
\usepackage[margin=1.9cm,top=1.6cm,bottom=1.8cm]{geometry}
\usepackage{microtype}
\usepackage{amsmath,amssymb}
\usepackage{booktabs}
\usepackage{tabularx}
\usepackage{array}
\usepackage{enumitem}
\usepackage{titlesec}
\usepackage{xcolor}
\usepackage[hidelinks]{hyperref}

% ------------------------------------------------------------------ layout
\setlength{\parindent}{0pt}
\setlength{\parskip}{3pt}
\titleformat{\section}{\normalsize\bfseries}{\thesection.}{0.5em}{}
\titlespacing*{\section}{0pt}{7pt}{3pt}
\setlist{nosep,leftmargin=1.2em,topsep=1pt}
\renewcommand{\arraystretch}{1.12}
\newcolumntype{Y}{>{\raggedright\arraybackslash}X}

\begin{document}

% ------------------------------------------------------------------ title block
\begin{center}
  {\Large\bfseries A Small Generative Model for Image Denoising}\\[3pt]
  {\large A Noise-Gated Conditional GAN with a NAF U-Net Generator}\\[8pt]
  \begin{tabular}{@{}c@{\hspace{3.5em}}c@{}}
    \textbf{Sagnik Chandra}        & \textbf{Shrusti Vipulbhai Jain} \\
    Roll No.\ M25CSE027            & Roll No.\ M25CSE030 \\
  \end{tabular}\\[6pt]
  \begin{tabular}{@{}r@{\hspace{0.6em}}l@{}}
    \textbf{Course:}  & Foundation Models and Generative AI \\
    \textbf{Faculty:} & Angshuman Paul \\
  \end{tabular}
\end{center}
\vspace{-6pt}
\noindent\rule{\linewidth}{0.4pt}
\vspace{-8pt}

% ==================================================================
\section{Overview}
We submit a \textbf{conditional generative adversarial network (cGAN)} that maps a corrupted
$256\times256\times3$ image, together with injected random noise, to a restored
$256\times256\times3$ uint8 image in a \textbf{single forward pass}. Different noise draws give
different plausible restorations (diversity 0.77 grey levels: mean per-pixel standard deviation
of the uint8 outputs over 8 noise seeds), so the model samples from a conditional distribution
rather than returning one fixed estimate. On an input of shape $1\times3\times256\times256$ the
generator costs \textbf{4.98\,GFLOPs} with fvcore and \textbf{5.09\,GMac} with ptflops (both
tools count one multiply-add as one operation) and has \textbf{5.52\,M} parameters; the discriminator is used only during training and adds no inference cost. The model
is trained on 26 corruption families, applied singly and in random combinations, over natural,
low-light and medical images.

% ==================================================================
\section{Architecture}
\textbf{Generator (NAF U-Net with decoder-only gated noise).} The generator is a five-scale U-Net
built from NAFBlocks (Chen et al., ECCV 2022). A $3\times3$ convolution lifts the input to 24
channels. The encoder has four levels at 256, 128, 64 and 32\,px with
$[2,2,2,4]$ blocks and widths $24,48,96,192$; each level ends with a $2\times2$ stride-2
convolution that halves the resolution and doubles the channels. A bottleneck of 3 blocks
operates at 16\,px with 384 channels. The decoder mirrors the encoder with $[1,1,1,2]$ blocks
(32\,px\,$\rightarrow$\,256\,px); each level upsamples with a $1\times1$ convolution and pixel
shuffle and adds the encoder feature map of the same scale (additive skip connection). A final
$3\times3$ convolution predicts a residual that is added to the corrupted input.
\begin{itemize}
  \item \textbf{NAFBlock:} LayerNorm $\rightarrow$ $1\times1$ conv $\rightarrow$ $3\times3$
        depthwise conv $\rightarrow$ SimpleGate (channel split, element-wise product)
        $\rightarrow$ simplified channel attention $\rightarrow$ $1\times1$ conv, followed by a
        gated feed-forward branch; both branches are residual with learned zero-initialised
        scales. The block uses no non-linear activation functions, which keeps it cheap.
  \item \textbf{Noise injection:} after each of the 5 decoder blocks, the feature map $x$
        receives $x + s_c\cdot\sigma(\mathrm{conv}_{1\times1}(x))\cdot n$, where
        $n\sim\mathcal{N}(0,1)$ is one noise value per pixel, $s_c$ is a learned per-channel
        strength (initialised to 0.01) and $\sigma(\mathrm{conv}_{1\times1}(x))$ is a learned
        per-pixel gate. The encoder and bottleneck are noise-free, so the skip connections carry
        clean features and the image structure is not perturbed; the gate lets the network learn,
        per pixel, where variation is added, so that it can be kept to uncertain regions.
  \item \textbf{Inference:} a fresh noise draw is made inside every forward call, so the module
        measured for FLOPs is exactly the one that runs. The output is clamped to $[0,1]$, multiplied by 255, rounded and returned as uint8.
\end{itemize}
\textbf{Discriminator (training only).} A conditional U-Net discriminator with spectral
normalisation (as in Real-ESRGAN) gives a real/fake score at every location of its output map. It operates in the
wavelet domain: its 18-channel input at 128\,px is the three high-frequency Haar bands of the
candidate and of the corrupted image, so it judges noise, blur and compression residue (which live
in those bands) while colour and brightness are left to the reconstruction losses. Seeing the
corrupted image makes it judge whether the output is the right restoration of \emph{that} input.

% ==================================================================
\section{Training objective}
For a corrupted image $y$, clean target $x$ and generator output $\hat{x}=G(y,n)$:
\begin{equation*}
\begin{aligned}
  \mathcal{L}_G ={}& \underbrace{\mathcal{L}_{\text{Charb}} + 0.05\,\mathcal{L}_{\text{FFT}}
     + 0.334\,(1-\mathrm{SSIM}) + 0.1\,\mathcal{L}_{\text{DISTS}}
     + 0.05\,\mathcal{L}_{\text{wav}}}_{\text{reconstruction, on every output}} \\
  & + 0.00275\,\mathcal{L}_{\text{adv}}
     + 4.0\,\max\!\bigl(0,\;\tau-\tfrac{1}{N}\lVert \hat{x}_1-\hat{x}_2\rVert_1\bigr).
\end{aligned}
\end{equation*}
\begin{itemize}
  \item \textbf{Reconstruction terms, one per scored metric:} Charbonnier loss for PSNR; L1
        distance between Fourier spectra (FFT) against over-smoothing; $1-\mathrm{SSIM}$ with the
        same $7\times7$ uniform window as \texttt{skimage} for SSIM; the DISTS metric itself as a
        loss (switched on after 50\% of training, on 4 images per batch) for perceptual quality; and L1 on Haar sub-bands
        ($\mathcal{L}_{\text{wav}}$).
  \item \textbf{Adversarial term:} hinge GAN loss, weight 0.00275, ramped in over the first 5k
        steps. The discriminator is trained with the hinge loss plus a lazy R1 gradient penalty
        ($\gamma=1$, every 16 steps).
  \item \textbf{Diversity constraint:} 4 images per batch are restored with a second noise draw;
        if the two samples differ by less than $\tau = 1.6$ grey levels (mean absolute
        difference), a hinge penalty pushes them apart, and it stops acting once the target is
        met. This keeps the noise input in use, so the model stays generative.
  \item \textbf{Optimisation:} AdamW, learning rate $7.55\times10^{-4}$ (discriminator
        $0.5\times$), 1k warm-up steps, cosine decay to $10^{-6}$, generator weight decay $10^{-4}$,
        gradient clipping 1.0, mixed precision, EMA of the weights (decay 0.999; the EMA weights
        are submitted). 200{,}000 steps of 16 images ($256\times256$) plus the 4 extra draws,
        trained from scratch in about 25\,h on one 24\,GB GPU.
  \item \textbf{Hyperparameter selection:} learning rate, discriminator learning-rate ratio,
        adversarial weight, diversity target and SSIM weight were chosen with Optuna (TPE
        sampler, successive-halving pruning, 38 short trials) under the constraint that diversity
        stays above 0.6 grey levels.
\end{itemize}

% ==================================================================
\section{Degradations and training data}
Every training pair is generated on the fly: a clean image is corrupted by a random chain of
\textbf{0 to 4} families (probabilities 12\%, 30\%, 36\%, 17\%, 5\%; the 12\% clean inputs teach
the model to leave good images untouched). Each family gets a continuous severity in $[0,1]$ drawn from
$\mathrm{Beta}(1,\,1.8+0.6(k-1))$ for a chain of $k$ families, so longer chains use milder
individual corruptions. Chains are applied in physical camera order (optics $\rightarrow$ blur
$\rightarrow$ noise $\rightarrow$ compression) or, for 50\% of images, in random order. A chain contains at most one family of each type (blur, codec, impulse, \dots) except noise,
of which up to two may be combined; the brief's nine families carry 65\% of the sampling
weight. Every corrupted image is quantised to uint8, as at test time.

\begin{center}
\small
\begin{tabularx}{\linewidth}{@{}p{3.6cm}Y@{}}
\toprule
\textbf{Group} & \textbf{Families} \\
\midrule
Listed in the brief (9) & Gaussian, Poisson, salt-and-pepper and speckle noise; Gaussian, motion
  (straight and curved trajectories) and defocus blur; JPEG; chromatic aberration \\
Additional (13) & spatially correlated, Rician, stripe/banding, sensor (signal-dependent) and
  periodic noise; random-valued impulses; resize blur; pixelation; ringing (over-sharpening halos); WebP; double JPEG;
  posterisation; small artifacts (scratches, dead pixels, small blocks) \\
Domain-specific (4) & demosaicing artifacts and chroma noise (cameras); MRI ghosting; CT ring
  artifacts \\
\bottomrule
\end{tabularx}
\end{center}

\textbf{Data.} Sources are sampled by weight rather than concatenated, so small domains are not
drowned out: \emph{natural} ($\approx$75\% of samples: DIV2K, Flickr2K, BSDS500, COCO), \emph{medical}
($\approx$22\%: NIH ChestX-ray X-rays, IXI brain MRI slices, MedMNIST CT and ultrasound) and
\emph{low-light} ($\approx$3\% real LOL images, plus 10\% of the natural images darkened
synthetically before corruption). Images
are resized so that their short side is close to 256\,px and randomly cropped to
$256\times256$, with flips, rotations and 4\% greyscale conversion. All are public research
datasets.

% ==================================================================
\section{Results}
Evaluated on our fixed held-out benchmark (396 images: 33 single, combined and unseen corruption
cases over natural, low-light and medical images), with \texttt{skimage} PSNR/SSIM and
\texttt{DISTS-pytorch} on uint8 images, as specified in the brief.
\begin{center}
\small
\begin{tabular}{@{}lccccc@{}}
\toprule
Model & PSNR (dB) $\uparrow$ & SSIM $\uparrow$ & DISTS $\downarrow$ & GFLOPs (fvcore) $\downarrow$ &
  Diversity$^{*}$ \\
\midrule
\textbf{Submitted model} & \textbf{32.74} & \textbf{0.8740} & \textbf{0.1371}
  & \textbf{4.98} & 0.77 \\
\bottomrule
\end{tabular}\\[2pt]
{\footnotesize $^{*}$Mean per-pixel standard deviation (grey levels, 0--255) of the uint8 outputs
over 8 noise seeds, on 96 benchmark images.}
\end{center}

\end{document}
